\section{Self-Improving Agents: un ciclo de retroalimentación no equivale a evolución autónoma}

La preference optimization actualiza la policy durante el entrenamiento; un agent, en cambio, percibe, planea, llama herramientas, actúa y recibe retroalimentación en un ciclo dentro del entorno. El curso divide la self-improvement en refinement, reflection, ReAct y tree search. Aquí, la improvement suele ser una mejora de la conducta en una sola tarea o una sola sesión; no debe reinterpretarse automáticamente como que el modelo aprendió de forma permanente una capacidad nueva.

\subsection{Definición mínima de un agent}

Los métodos de la sección anterior agregan razonamiento o actualizan parámetros, sobre todo, dentro de una sola solicitud; un agent, en cambio, conecta el modelo a un state y un environment que persisten. El error más común al definir un agent es dibujar solo una flecha circular sin indicar de dónde vienen las observaciones, qué permisos tienen las acciones, si la memory es confiable y cuándo se detiene. La ecuación de estado siguiente hace explícitos estos objetos, de modo que refinement y tree search tengan después un punto de aplicación claro.

Un agent puede abstraerse como
\[
s_{t+1}=F(s_t,o_t,m_t),\qquad a_t=\pi_{\theta}(s_t,\mathcal{T}),
\]
donde $o_t$ es la observación, $m_t$ es la memory, $\mathcal{T}$ es el conjunto de herramientas y $a_t$ es la acción; el entorno devuelve entonces la siguiente observación. El modelo base es solo una parte de $\pi_{\theta}$; los permisos, el estado, el schema de las herramientas, la recuperación ante errores y las condiciones de paro también determinan la conducta del sistema.


\lecturefigure{slide-085.jpg}{What is an AI Agent: entorno, modelo, objetivo, herramientas, memoria y acciones}{CS25 V5 official deck, p.~85}

\teachervoice{00:34:20--00:35:20, Chelsea define un agent como un sistema capaz de percibir el entorno, tomar decisiones y actuar, y enumera memory, tools y feedback. Nota de clase: envolver una llamada al modelo en una función no basta para obtener agent reliability.}

\subsection{Refinement y reflection}

Refinement hace que el modelo genere una respuesta, encuentre sus debilidades y la revise; reflection pone más énfasis en extraer de cada failure experiencia reutilizable y escribirla en la memory. Ambas pueden caer en un confirmation loop: si el mismo modelo responde y revisa, un supuesto erróneo puede reforzarse una y otra vez.


\lecturefigure{slide-086.jpg}{Self-improvement: refinement, reflection, exploración y razonamiento por múltiples rutas}{CS25 V5 official deck, p.~86}


\lecturefigure{slide-087.jpg}{Refinement: generación, evaluación, retroalimentación y revisión iterativa}{CS25 V5 official deck, p.~87}


\lecturefigure{slide-088.jpg}{Reflection: extraer experiencia de los fallos y actualizar la estrategia posterior}{CS25 V5 official deck, p.~88}

\begin{warningbox}{La self-critique necesita un ancla externa}
Sin tool outcome, test, human check o un verifier independiente, «ya lo revisé» es solo otro fragmento generado. En ingeniería, la critique debe convertirse en aserciones ejecutables que luego verifique el entorno; lo que no pueda verificarse debe conservar su uncertainty.
\end{warningbox}

\subsection{ReAct y LATS: conectar el reasoning con action/search}

ReAct alterna entre thought y action, de modo que el modelo actualiza su plan según la observation. LATS extiende la ruta única a múltiples action trajectories y agrega la feedback con una selección MCTS-like; aumenta la exploración, pero multiplica el costo de las llamadas a herramientas, de la gestión del estado y de las ramas erróneas.


\lecturefigure{slide-089.jpg}{ReAct: alternancia de reasoning, action y observation}{CS25 V5 official deck, p.~89}


\lecturefigure{slide-090.jpg}{LATS: agent search por múltiples rutas, agregación de retroalimentación y búsqueda en árbol}{CS25 V5 official deck, p.~90}

\begin{warningbox}{Protocolo de confiabilidad para la agent improvement}
Cada iteration debe guardar el state de entrada, el plan del modelo, la action real, el tool output, la verifier evidence y la memory write; de lo contrario, cuando el sistema falle no podrá distinguirse si fue un error de planeación, de herramienta, de permisos o una contaminación de la memoria. Para refinement, mide la exactitud después de la revisión y los errores nuevos que introduce; para reflection, mide si la memoria realmente mejora tareas futuras similares sin propagar errores; para ReAct, mide los parámetros de las herramientas, el observation parsing y las condiciones de paro; para LATS, mide el beneficio marginal de aumentar el presupuesto de búsqueda y el evaluator bias. En producción también se necesitan action idempotency, transaction rollback, secret isolation y una ruta de escalamiento a personas. Lo que se llama self-improving solo es una capacidad de ingeniería si reduce los fallos de forma sostenida con evidence auditable, no si el modelo genera más texto de autoevaluación.
\end{warningbox}

\begin{knowledgebox}{Lectura de la figura: las cuatro formas de self-improvement modifican estados distintos}
Refinement reescribe la answer actual; reflection escribe la experiencia de los fallos en la memory; ReAct actualiza la trajectory actual con la observation; LATS mantiene en el search tree varias candidatas junto con su valor. La estructura de árbol de la Slide 90 respalda «explorar más rutas», pero no demuestra que la puntuación de la búsqueda sea confiable. Al leer la figura, marca primero state, action, feedback y persistent memory, y luego pregúntate si el entorno puede verificar los errores. Si la feedback la sigue generando el mismo modelo sin ningún ancla, el ciclo puede limitarse a redactar el error original de forma más completa.
\end{knowledgebox}

\teachervoice{00:35:20--00:39:10, en clase se distinguen sucesivamente refinement, reflection, ReAct y LATS. Experiencia práctica: la diferencia entre estos mecanismos está en dónde se escribe la feedback: en el texto actual, en la memory de largo plazo, en el state del entorno o en el search tree.}

\begin{lstlisting}[caption={Agent loop con verificación y límites de permisos}]
state = initialize(goal, permissions, budget)
while not done(state):
    proposal = model.plan(state)
    action = policy_gate(proposal, permissions)
    observation = execute_in_sandbox(action)
    evidence = verify_outcome(observation)
    state = update_memory(state, action, evidence)
    if repeated_failure(state): escalate_or_stop(state)
\end{lstlisting}

\subsection{Resumen de la sección}

La capacidad de un agent surge del ciclo cerrado entre el modelo y el entorno. Refinement reescribe respuestas, reflection escribe experiencia, ReAct conecta herramientas y LATS agrega búsqueda por múltiples rutas; una self-improvement realmente confiable depende de outcomes verificables, control de permisos, presupuesto y paro ante fallos, no de un diálogo consigo mismo sin fin.
